% !TeX root = ../../main.tex
\subsection{割合の問題}

    \begin{techniquebox}[tech:04]{\textbf{割合を扱うときの鉄則}}
        \textbf{割合は\fitblankbf[分数]{}や\fitblankbf[小数]{}に直して
        \fitblankbf[かけ算]{}で扱う！}
    \end{techniquebox}

    割合の問題のうち，代表的なものに，\uwave{「生徒数の問題」}と
    \uwave{「食塩水の問題」}がある.それぞれについて，深掘りしていこう.

    \subsubsection{生徒数の問題}
    \vspace{1em}

        \begin{techniquebox}[tech:05]{\textbf{生徒数の問題}}
            \textbf{(男子の生徒数) + (女子の生徒数) = (全校生徒数)} という考え方を使う!
        \end{techniquebox}

        Technique \ref{tech:05} は最低限使えて欲しい考え方ではあるが，余裕があれば，
        \uwave{生徒数の増減に着目した}考え方も使えて欲しい.

        \[\textbf{(男子の増減) + (女子の増減) = (その年の生徒数の増減)}\]

        \vspace{1.5em}
        という式を立てることができると計算が楽になる.

% 例題３
        \begin{example}[【例題 3】]
            ある中学校の昨年の入学者数は360人であった.
            今年は昨年より男子が3\%減少し，女子が5\%増加して，
            全体では2人増加したという.昨年の男子，女子の入学者数をそれぞれ求めなさい.
        \end{example}

        \begin{proof}\mbox{}\\
            \[（昨年の男子の入学者数）+（昨年の女子の入学者数）=（昨年の入学者数）\]
            \[（今年の男子の入学者数）+（今年の女子の入学者数）=（今年の入学者数）\]

            という方針にしたがって，方程式を立てる.

            \begin{hiddenanswer}
                昨年の男子の入学者数を $x$ 人, 女子の入学者数を $y$ 人とすると，

                \begin{align*}
                    \left\{
                    \begin{array}{@{} r @{\quad+\quad} l @{\quad=\quad} r @{}}
                        x\phantom{\times \dfrac{97}{100}} & y\phantom{\times \dfrac{105}{100}} & 360 \\
                        x \times \dfrac{97}{100} & y \times \dfrac{105}{100} & 362
                    \end{array}
                    \right.
                \end{align*}

                これを解いて，$x = 200, y = 160$ となる.

                \finalanswer{答え：昨年の男子200人,女子160人}
            \end{hiddenanswer}
        \end{proof}

    \subsubsection{食塩水の問題}
    \vspace{1em}

        \begin{techniquebox}[tech:06]{\textbf{食塩水の問題}}
            \textbf{(混ぜる前の塩の量) = (混ぜた後の塩の量)} という考え方を使う!
        \end{techniquebox}

        塩の量は，(食塩水の量) $\times$ (塩の濃度) で求められるので，Technique \ref{tech:04} より，

        \[\textbf{(塩の量) = (食塩水の量)} \times \left(\dfrac{\text{\textbf{濃度}}}{\bm{100}}\right)\]

% 例題４
        \begin{example}[【例題 4】]
            5\%の食塩水と8\%の食塩水を混ぜて，6\%の食塩水を300gつくりたい.
            2種類の食塩水をそれぞれ何gずつ混ぜればよいか求めなさい.
        \end{example}

        \begin{proof}\mbox{}\\
            \begin{align*}
                （5\%の食塩水の量）+（8\%の食塩水の量） &=（6\%の食塩水の量） \\
                （5\%の食塩水の塩の量）+（8\%の食塩水の塩の量） &=（6\%の食塩水の塩の量）
            \end{align*}

            という方針にしたがって，方程式を立てる.

            \begin{hiddenanswer}
                5\%の食塩水の量を $x$ g, 8\%の食塩水の量を $y$ gとすると，

                \begin{align*}
                    \left\{
                    \begin{array}{@{} r @{\quad+\quad} l @{\quad=\quad} r @{}}
                        x\phantom{\times \dfrac{5}{100}} & y\phantom{\times \dfrac{8}{100}} & 300 \\
                        x \times \dfrac{5}{100} & y \times \dfrac{8}{100} & 300 \times \dfrac{6}{100}
                    \end{array}
                    \right.
                \end{align*}

                これを解いて，$x = 200, y = 100$ となる.

                \finalanswer{答え：5\%の食塩水200g, 8\%の食塩水100g}
            \end{hiddenanswer}
        \end{proof}
\newpage
